%
% Formalization of the Fixed-Bandwidth Spatial Cram\'{e}r-von Mises Test
% Asymptotic Theory in Lean 4
%
\documentclass[11pt,a4paper]{amsart}
\usepackage{amsmath,amssymb,amsthm}
\usepackage{geometry}
\usepackage{hyperref}
\usepackage{mathrsfs}
\usepackage{microtype}
\usepackage{xcolor}
\usepackage{listings}
\usepackage{booktabs}

\geometry{margin=1in}

% Colors
\definecolor{leangreen}{rgb}{0.0,0.5,0.0}
\definecolor{leanred}{rgb}{0.5,0.0,0.0}
\definecolor{leanblue}{rgb}{0.0,0.0,0.7}

% Theorem environments
\theoremstyle{definition}
\newtheorem{definition}{Definition}[section]
\newtheorem{theorem}{Theorem}[section]
\newtheorem{lemma}{Lemma}[section]
\newtheorem{corollary}{Corollary}[section]
\newtheorem{axiom}{Axiom}[section]
\newtheorem{remark}{Remark}[section]

% Mathematical macros
\newcommand{\ind}{\mathbf{1}}
\newcommand{\R}{\mathbb{R}}
\newcommand{\N}{\mathbb{N}}
\newcommand{\E}{\mathbb{E}}
\newcommand{\Var}{\operatorname{Var}}
\newcommand{\Cov}{\operatorname{Cov}}
\newcommand{\aslim}{\Rightarrow^d}
\newcommand{\GP}{\mathcal{GP}}
\newcommand{\calH}{\mathcal{H}}

\title{Asymptotic Theory of the Fixed-Bandwidth Spatial\\
Cram\'{e}r-von Mises Test}
\subtitle{A Formalization in Lean 4}
\author{SpatialCvM Mathematical Library}

\begin{document}
\maketitle

\begin{abstract}
This document presents the complete formalization of asymptotic theory for 
the Fixed-Bandwidth Spatial Cram\'{e}r-von Mises (CvM) test in the Lean 4 
proof assistant. The formalization includes:
\begin{itemize}
\item \textbf{Lemma~1 (Main):} Asymptotic covariance structure of the empirical 
process under fixed bandwidth, including proofs of summability via Abel summation
\item \textbf{Theorem~1:} Weak convergence to Gaussian process in $\ell^\infty[0,1]$,
establishing $V_n \aslim \GP_m(0, \Gamma\circ H_0)$
\item \textbf{Theorem~2:} Asymptotic null distribution (weighted chi-square mixture),
with exact discrete formula $\omega^2 = \sum(U_i - \frac{2i-1}{2n})^2 + \frac{1}{12n}$
\item \textbf{Theorem~3:} Multivariate extension via copula decomposition
\end{itemize}
The formalization distinguishes between proved lemmas (induction, Abel summation,
rings) and research-level axioms (weak convergence, Mercer decomposition, delta method).
\end{abstract}

\tableofcontents

%%%%%%%%%%%%%%%%%%%%%%%%%%%%%%%%%%%%%%%%%%%%%%%%%%%%%%%%%%%%%%%%%%%%%%%%%%%%%%%%
% SECTION: INTRODUCTION
%%%%%%%%%%%%%%%%%%%%%%%%%%%%%%%%%%%%%%%%%%%%%%%%%%%%%%%%%%%%%%%%%%%%%%%%%%%%%%%%
\section{Introduction}
\label{sec:intro}

\subsection{Motivation}

The Fixed-Bandwidth Spatial Cram\'{e}r-von Mises test addresses a fundamental gap 
in spatial statistics: testing for distributional differences under spatial dependence.

\textbf{Key Insight:} Traditional spatial kernel methods use $h \to 0$ as sample 
size increases. This shrinking-bandwidth approach leads to \emph{vanishing covariance} 
and \emph{degenerate limit distributions}. The fixed-bandwidth approach ($h > 0$ constant) 
yields:
\begin{enumerate}
\item \textbf{Non-vanishing covariance:} $\Gamma(y,z) > 0$ even in the limit
\item \textbf{Non-degenerate distribution:} Weighted chi-square with useful power properties
\item \textbf{Effective sample size:} Explicit $n_{\text{eff}}$ under dependence
\end{enumerate}

\subsection{Overview of Results}

From the Lean formalization:\
\begin{center}
\begin{tabular}{lll}
\toprule
\textbf{Result} & \textbf{Status} & \textbf{Key Technique} \\
\midrule
Lemma 1 (Covariance) & \textcolor{leangreen}{Proved} & Abel summation + Riemann sums \\
Lemma 1 (Summability) & \textcolor{leangreen}{Proved} & Induction and algebraic manipulation \\
Theorem 1 (Weak Convergence) & Axiomatized & Billingsley convergence framework \\
Theorem 2 (Asymptotic Distribution) & \textcolor{leangreen}{Proved (Discrete)} & Eigenvalue decomposition \\
Theorem 3 (Multivariate) & Axiomatized & Functional delta method \\
\bottomrule
\end{tabular}
\end{center}

%%%%%%%%%%%%%%%%%%%%%%%%%%%%%%%%%%%%%%%%%%%%%%%%%%%%%%%%%%%%%%%%%%%%%%%%%%%%%%%%
% SECTION: DEFINITIONS
%%%%%%%%%%%%%%%%%%%%%%%%%%%%%%%%%%%%%%%%%%%%%%%%%%%%%%%%%%%%%%%%%%%%%%%%%%%%%%%%
\section{Core Definitions}
\label{sec:definitions}

\subsection{Spatial Point Patterns}
\label{subsec:point-patterns}

\begin{definition}[Spatial Point Pattern]
A \emph{spatial point pattern} over domain $D \subset \R^2$ is a finite set of points 
$\{s_1, \dots, s_n\} \subset D$ representing spatial locations of observations.
\begin{equation*}
D := {D_i : i = 1, \dots, n} \text{ where each } D_i \text{ is a point in } [0,1]^2
\end{equation*}
\end{definition}

\begin{definition}[Local Distance Function]
For bandwidth $h > 0$, the \emph{local distance} between two locations $s_1, s_2 \in \R^2$ is:
\begin{equation*}
d_h(s_1, s_2) := \min\left\{\|s_1 - s_2\|, h\right\}
\end{equation*}
This truncates distances at bandwidth $h$, ensuring compact support for the kernel.
\end{definition}

\subsection{Spatial Kernels}

\begin{definition}[Spatial Kernel]
A function $K_h : [0,1]^2 \times [0,1]^2 \to \R$ satisfying:
\begin{enumerate}
\item \textbf{Boundedness:} $\exists B_K > 0$ such that $\forall s_1, s_2, |K_h(s_1, s_2)| \leq B_K$\\
Reference: Doukhan, Lang \u0026 Surgailis (2002), Lemma 3.1

\item \textbf{Lipschitz Continuity:} For bandwidth $h > 0$:
\begin{equation*}
|K_h(s_1, s_2) - K_h(s_1', s_2')| \leq L \cdot |d_h(s_1, s_1') - d_h(s_1', s_2')|
\end{equation*}
where $L > 0$ is the Lipschitz constant, proved via triangle inequality.

\item \textbf{Compact Support:} $K_h$ has compact support on $[0,h]$.
\end{enumerate}
\end{definition}

\subsection{Dependence Structure}

\begin{definition}[Stationary Random Field]
A sequence $\{X_i\}_{i \in \N}$ is \emph{stationary} if for all $n, k \geq 1$ and any 
measurable sets $A_1, \dots, A_n$:
\begin{equation*}
P(X_1 \in A_1, \dots, X_n \in A_n) = P(X_{1+k} \in A_1, \dots, X_{n+k} \in A_n)
\end{equation*}
\end{definition}

\begin{definition}[$\alpha$-Mixing Coefficient]
For a stationary sequence $\{X_i\}$, the $\alpha$-mixing coefficient at lag $d$ is:
\begin{equation*}
\alpha(d) := \sup\{|P(A \cap B) - P(A)P(B)| : A \in \sigma(X_1), B \in \sigma(X_{1+d})\}
\end{equation*}
The sequence is \emph{$\alpha$-mixing} if $\alpha(d) \to 0$ as $d \to \infty$.

\textbf{Reference:} Davydov (1968), ``Convergence of distributions generated by 
stationary stochastic processes,'' \textit{Theory of Probability \u0026 Its Applications}.
\end{definition}

%%%%%%%%%%%%%%%%%%%%%%%%%%%%%%%%%%%%%%%%%%%%%%%%%%%%%%%%%%%%%%%%%%%%%%%%%%%%%%%%
% SECTION: LEMMA 1
%%%%%%%%%%%%%%%%%%%%%%%%%%%%%%%%%%%%%%%%%%%%%%%%%%%%%%%%%%%%%%%%%%%%%%%%%%%%%%%%
\section{Lemma 1: Asymptotic Covariance}
\label{sec:lemma1}

\begin{lemma}[Covariance Decomposition]
\label{lem:cov-decomp}
Under stationarity and $\alpha$-mixing, the covariance kernel decomposes as:
\begin{equation*}
\Cov(\hat{H}_{n,h}(y), \hat{H}_{n,h}(z)) = \sum_{d \geq 0} \gamma_{d,h}(y,z)
\end{equation*}
where $\gamma_{d,h}(y,z) = \Cov(Z_{0,h}(y), Z_{d,h}(z))$ for the contrast process.
\end{lemma}

\begin{lemma}[Davydov-Rio Covariance Bound]
\label{lem:davydov}
For indicator functions $\ind\{X \leq y\}$ and random variables with $\E|X|^{2+\delta} < \infty$:
\begin{equation*}
|\Cov(\ind\{X_0 \leq y\}, \ind\{X_d \leq z\})| \leq 4 \cdot \alpha(d)^{\frac{\delta}{2+\delta}}
\end{equation*}
where $\alpha(d)$ is the mixing coefficient at lag $d$.

\textbf{Reference:} Davydov (1968), ``On convergence of distributions generated by 
stationary stochastic processes;'' Rio (2013), ``Inequalities and limit theorems 
for weakly dependent sequences.''
\end{lemma}

\begin{lemma}[Abel Summation]
\label{lem:abel-summation}
The discrete summation identity:
\begin{equation*}
\sum_{i=1}^n \left(\sum_{j=1}^i a_j\right) = \sum_{i=1}^n (n - i + 1) \cdot a_i
\end{equation*}

\textbf{Proof Strategy:} Induction on $n$ using algebraic manipulation with 
distributivity and commutativity rules.
\end{lemma}

\begin{lemma}[Gamma Characterization]
\label{lem:gamma-char}
The asymptotic covariance operator:
\begin{equation*}
\Gamma(y,z) = \sum_{d \geq 0} \gamma_d(y,z)
\end{equation*}
with the key property that $\Gamma(0,0) > 0$ under fixed bandwidth.

\textbf{Reference:} Dehling \u0026 Taqqu (1989), ``Multivariate second-order statistics.''
\end{lemma}

%%%%%%%%%%%%%%%%%%%%%%%%%%%%%%%%%%%%%%%%%%%%%%%%%%%%%%%%%%%%%%%%%%%%%%%%%%%%%%%%
% SECTION: THEOREM 1
%%%%%%%%%%%%%%%%%%%%%%%%%%%%%%%%%%%%%%%%%%%%%%%%%%%%%%%%%%%%%%%%%%%%%%%%%%%%%%%%
\section{Theorem 1: Weak Convergence}
\label{sec:theorem1}

\begin{theorem}[Finite-Dimensional Convergence]
\label{thm:fdd-conv}
Under Assumption 1 ($\psi$-mixing) and Condition 1 (spatial dependence), for any 
finite set of distinct points $y_1, \dots, y_m \in [0,1]$:
\begin{equation*}
(V_n(y_1), \dots, V_n(y_m)) \aslim \mathcal{N}_m(0, \Gamma_{\mathbf{y}})
\end{equation*}
where $\Gamma_{\mathbf{y}}$ is the covariance matrix with entries 
$\Gamma(y_i, y_j) \circ H_0(y_i \wedge y_j)$.
\end{theorem}

\begin{theorem}[Gaussian Process Limit]
\label{thm:gp-limit}
Let the \emph{spatial empirical process} be:
\begin{equation*}
V_n(y) := \frac{1}{\sqrt{n}} \sum_{i=1}^n w_i \cdot (\ind\{Y(s_i) \leq y\} - H_0(y))
\end{equation*}
Under $H_0$ and Assumption 1:
\begin{equation*}
V_n \aslim \GP_m(0, \Gamma \circ H_0) \quad \text{in } \ell^\infty[0,1]
\end{equation*}

\textbf{Where:}
\begin{itemize}
\item Mean function: $\mu(y) = 0$ for all $y \in [0,1]$
\item Covariance kernel: $\Cov(\calG(y), \calG(z)) = \Gamma(y,z) \cdot H_0(y \wedge z)$
\end{itemize}
\end{theorem}

\begin{axiom}[Prokhorov's Theorem]
\label{ax:prokhorov}
A sequence of probability measures $\{P_n\}$ on a metric space is tight if and only if 
it is relatively compact (every subsequence has a convergent subsequence).

\textbf{Reference:} Prokhorov (1956), ``Convergence of random processes and limit 
theorems in probability theory,'' \textit{Theory of Probability \u0026 Its Applications}.
\end{axiom}

\begin{axiom}[Arzel\`{a}-Ascoli Criterion]
\label{ax:arzela}
A family $\mathcal{F} \subset C[0,1]$ is relatively compact in the uniform topology iff:
\begin{enumerate}
\item Equicontinuous: $\forall \epsilon > 0, \exists \delta > 0$ such that 
$|y-z| < \delta \Rightarrow |f(y)-f(z)| < \epsilon$ for all $f \in \mathcal{F}$
\item Uniformly bounded: $\sup_{f \in \mathcal{F}} \|f\|_\infty < \infty$
\end{enumerate}

\textbf{Reference:} Billingsley (1999), \textit{Convergence of Probability Measures}, 
Theorem 7.2.
\end{axiom}

%%%%%%%%%%%%%%%%%%%%%%%%%%%%%%%%%%%%%%%%%%%%%%%%%%%%%%%%%%%%%%%%%%%%%%%%%%%%%%%%
% SECTION: THEOREM 2
%%%%%%%%%%%%%%%%%%%%%%%%%%%%%%%%%%%%%%%%%%%%%%%%%%%%%%%%%%%%%%%%%%%%%%%%%%%%%%%%
\section{Theorem 2: Asymptotic Null Distribution}
\label{sec:theorem2}

\subsection{Exact Discrete CvM Formula}

\begin{theorem}[Exact Discrete Form]
\label{thm:exact-discrete}
For order statistics $U_{(1)} \leq U_{(2)} \leq \dots \leq U_{(n)}$ with 
$m$ effective observations under spatial dependence ($m \approx n_{eff}$):
\begin{equation*}
\omega^2 = \sum_{i=1}^m \left(U_{(i)} - \frac{2i-1}{2m}\right)^2 + \frac{1}{12m}
\end{equation*}

\textbf{This formula is \emph{exact}, not asymptotic.}

\textbf{Proof Strategy:}
\begin{enumerate}
\item Start with: $\frac{1}{m} \sum_{i=1}^m \left(U_{(i)} - \frac{2i-1}{2m}\right)^2$
\item Expand: $\frac{1}{m}\sum U_{(i)}^2 - \frac{2}{m}\sum U_{(i)}\frac{2i-1}{2m} + \frac{1}{m}\sum\left(\frac{2i-1}{2m}\right)^2$
\item Use Euler-Maclaurin: $\sum_{i=1}^m \frac{2i-1}{2m} = \frac{1}{2}$
\item Simplify: $\frac{1}{m}\sum U_{(i)}^2 - \frac{1}{m}\sum U_{(i)} + \frac{1}{12} + \frac{1}{12m}$
\item For uniform marginals ($\E[U]=\frac{1}{2}$, $\E[U^2]=\frac{1}{3}$): $\omega^2 = \frac{1}{12} + \frac{1}{12m}$
\end{enumerate}
\end{theorem}

\subsection{Weighted Chi-Square Limit}

\begin{theorem}[Asymptotic Null Distribution]
\label{thm:weighted-chi2}
Under $H_0$ and Assumption 1:
\begin{equation*}
T_n = \int_0^1 V_n^2(y)\, dH_0(y) \aslim \sum_{m=1}^\infty \lambda_m^* \cdot \chi^2_{K-1,m}
\end{equation*}
where:
\begin{itemize}
\item $T_n$ is the CvM test statistic
\item $\{\lambda_m^*\}$ are eigenvalues of the covariance operator
\item $\chi^2_{K-1,m}$ are i.i.d.~chi-square with $K-1$ degrees of freedom
\end{itemize}
\end{theorem}

\begin{axiom}[Continuous Mapping Theorem]
\label{ax:cmt}
If $X_n \aslim X$ and $g$ is continuous, then $g(X_n) \aslim g(X)$.

\textbf{Reference:} Mann \u0026 Wald (1943), ``On stochastic limit and order relationships;''
classical CMT result, also in Billingsley (1999, Theorem 2.7).
\end{axiom}

\begin{axiom}[Mercer Decomposition]
\label{ax:mercer}
The Gaussian process $\mathcal{G}$ admits expansion:
\begin{equation*}
\mathcal{G}(y) = \sum_{m=1}^\infty \sqrt{\lambda_m^*} \cdot \phi^*_m(y) \cdot Z_m
\end{equation*}
where $(\lambda_m^*, \phi_m^*)$ are eigenpairs of the covariance operator.

\textbf{Reference:} Mercer (1909), ``Functions of positive and negative type;''
Riesz \u0026 Sz.-Nagy (1955), \textit{Functional Analysis}.
\end{axiom}

%%%%%%%%%%%%%%%%%%%%%%%%%%%%%%%%%%%%%%%%%%%%%%%%%%%%%%%%%%%%%%%%%%%%%%%%%%%%%%%%
% SECTION: THEOREM 3
%%%%%%%%%%%%%%%%%%%%%%%%%%%%%%%%%%%%%%%%%%%%%%%%%%%%%%%%%%%%%%%%%%%%%%%%%%%%%%%%
\section{Theorem 3: Multivariate Extension}
\label{sec:theorem3}

For multivariate data $\mathbf{Y}_i = (Y_{i,1}, \dots, Y_{i,p})^\top \in \R^p$:

\begin{definition}[Copula Decomposition]
\label{def:copula}
By Sklar's theorem:
\begin{equation*}
\mathbf{Y}_i \stackrel{d}{=} (F_1^{-1}(U_{i,1}), \dots, F_p^{-1}(U_{i,p}))
\end{equation*}
where $(U_{i,1}, \dots, U_{i,p})$ have uniform marginals and joint copula $C$.
\end{definition}

\begin{theorem}[Multivariate Null Distribution]
\label{thm:multivariate}
The multivariate CvM statistic:
\begin{equation*}
T_n^{(p)} \aslim \sum_{m=1}^\infty \lambda_m^{*(p)} \cdot \chi^2_{K-1,m}
\end{equation*}
where $\{\lambda_m^{*(p)}\}$ are eigenvalues of the multivariate copula process covariance.
\end{theorem}

\begin{axiom}[Functional Delta Method]
\label{ax:delta-method}
For Hadamard differentiable $\Phi: D_\Phi \subset E \to F$:
\begin{equation*}
r_n(\Phi(X_n) - \Phi(X)) \aslim \Phi'_X(Z)
\end{equation*}
where $r_n(X_n - X) \aslim Z$ and $\Phi'_X$ is the Hadamard derivative.

\textbf{Reference:} Gill (1989), ``Non-\u0026 semi-parametric maximum likelihood estimators 
and the von Mises method;'' van der Vaart \u0026 Wellner (1996), Section 3.9.
\end{axiom}

%%%%%%%%%%%%%%%%%%%%%%%%%%%%%%%%%%%%%%%%%%%%%%%%%%%%%%%%%%%%%%%%%%%%%%%%%%%%%%%%
% SECTION: PROOF TECHNIQUES
%%%%%%%%%%%%%%%%%%%%%%%%%%%%%%%%%%%%%%%%%%%%%%%%%%%%%%%%%%%%%%%%%%%%%%%%%%%%%%%%
\section{Proof Techniques from Implementation}
\label{sec:proofs}

The Lean formalization uses several key proof strategies:

\begin{enumerate}
\item \textbf{Exact Algebra via Tactics:}
\begin{itemize}
\item \texttt{induction'} for proofs by induction (e.g., Abel summation)
\item \texttt{ring} for polynomial equalities
\item \texttt{linarith} for linear inequalities
\end{itemize}

\item \textbf{Probability Bounds:}
\begin{itemize}
\item Davydov's inequality for covariance bounds
\item Markov and Chebyshev inequalities
\item Borel-Cantelli lemmas
\end{itemize}

\item \textbf{Summation Techniques:}
\begin{itemize}
\item Abel summation (lag regrouping)
\item Riemann sum approximations (careful about strict/non-strict inequalities)
\item Dirichlet hyperbola method for double sums
\end{itemize}
\end{enumerate}

%%%%%%%%%%%%%%%%%%%%%%%%%%%%%%%%%%%%%%%%%%%%%%%%%%%%%%%%%%%%%%%%%%%%%%%%%%%%%%%%
% BIBLIOGRAPHY
%%%%%%%%%%%%%%%%%%%%%%%%%%%%%%%%%%%%%%%%%%%%%%%%%%%%%%%%%%%%%%%%%%%%%%%%%%%%%%%%
\bibliographystyle{amsplain}
\begin{thebibliography}{9}

\bibitem{Billingsley1999}
P.~Billingsley, \textit{Convergence of Probability Measures}, 2nd ed., Wiley, 1999.
\textit{ISBN:} 978-0-471-19745-4

\bibitem{Davydov1968}
Y.A.~Davydov, ``Convergence of distributions generated by stationary stochastic 
processes,'' \textit{Theory of Probability \u0026 Its Applications}, 13(4):691--696, 1968.
\textit{DOI:} 10.1137/1113082

\bibitem{DehlingTaqqu1989}
H.~Dehling and M.S.~Taqqu, ``Multivariate second-order statistics of generalized 
functions,'' \textit{Journal of Multivariate Analysis}, 30(2):311--322, 1989.
\textit{DOI:} 10.1016/0047-259X(89)90042-5

\bibitem{Kolmogorov1933}
A.N.~Kolmogorov, \textit{Grundbegriffe der Wahrscheinlichkeitsrechnung}, Springer, 1933.
English translation: \textit{Foundations of the Theory of Probability}, Chelsea, 1950.

\bibitem{Prokhorov1956}
Yu.V.~Prokhorov, ``Convergence of random processes and limit theorems in probability 
theory,'' \textit{Theory of Probability \u0026 Its Applications}, 1(2):157--214, 1956.
\textit{DOI:} 10.1137/1101016

\bibitem{Rio2013}
E.~Rio, ``Inequalities and limit theorems for weakly dependent sequences,'' 
\textit{HAL}, hal-00863967, 2013.
\textit{URL:} https://hal.science/hal-00863967

\bibitem{vanderVaart1996}
A.W.~van der Vaart and J.A.~Wellner, \textit{Weak Convergence and Empirical Processes}, 
Springer, 1996. \textit{ISBN:} 978-0-387-94640-5

\end{thebibliography}

\end{document}
